% Figure from MATH 452 notes, section 22. Compile with the notes preamble.
\begin{tikzpicture}[scale=1.0]
% left box: parameter domain
\node[draw, thick, figB!70, fill=figB!6, minimum height=1.5cm, align=center, font=\small, inner xsep=8pt] (Dbox) at (0,0) {parameter domain\\ $D \subseteq \mathbb{R}^k$, coords $(u_j)$};
% right box: ambient space
\node[draw, thick, figB!70, fill=figB!6, minimum height=1.5cm, align=center, font=\small, inner xsep=8pt] (Rbox) at (9.2,0) {ambient space\\ $\mathbb{R}^n$, coords $(x_i)$};
% forward arrow: points
\draw[->, very thick, figG] ($(Dbox.east)+(0.1,0.34)$) -- ($(Rbox.west)+(-0.1,0.34)$)
  node[midway, above, font=\scriptsize, figG!75!black] {$\Phi$: points travel forward};
% backward arrow: forms
\draw[<-, very thick, figR] ($(Dbox.east)+(0.1,-0.34)$) -- ($(Rbox.west)+(-0.1,-0.34)$)
  node[midway, below, font=\scriptsize, figR!85!black] {$\Phi^*$: forms travel backward};
% annotations under each box
\node[font=\scriptsize, align=center] at (0,-1.6) {$\Phi^*\omega$ lives here\\ (ready to integrate: $\int_D \Phi^*\omega$)};
\node[font=\scriptsize, align=center] at (9.2,-1.6) {$\omega$ lives here\\ (written in $dx_i$)};
\end{tikzpicture}
